\section*{Chapitre \ref{ch:in:structurer-sales-and-sales-trader} --- Structureur, vendeur et sales-trader}

\begin{solution}{exo:in:structurer-sales-and-sales-trader:1}
$(286\,264+99\,890)/486\,635=79.4~\%$.
\end{solution}

\begin{solution}{exo:in:structurer-sales-and-sales-trader:2}
Une marge de 1.83 pour 100~; un coupon de 6.97~\% par an n'en laisserait aucune.
\end{solution}

\begin{solution}{exo:in:structurer-sales-and-sales-trader:3}
Article 16(3)~: un processus d'approbation de chaque instrument financier avant sa commercialisation, précisant un marché cible identifié
de clients finaux, évaluant les risques qui le concernent et vérifiant la stratégie de distribution par rapport à lui, avec un réexamen régulier.
\end{solution}

\begin{solution}{exo:in:structurer-sales-and-sales-trader:4}
La note est une obligation plus des options que l'investisseur vend. À un taux plus élevé, la part obligataire coûte moins cher aujourd'hui pour un même
remboursement, ce qui libère de la valeur pour le coupon~; le forward de l'indice augmente aussi, ce qui rend plus probable le remboursement anticipé.
\end{solution}

\begin{solution}{exo:in:structurer-sales-and-sales-trader:5}
Il vend davantage de volatilité~: à 30~\% de volatilité, la probabilité que l'indice finisse sous 60~\% et qu'il perde son capital passe de 9.4~\% à
19.3~\% des trajectoires, et le coupon compense ce risque.
\end{solution}

\begin{solution}{exo:in:structurer-sales-and-sales-trader:6}
Le desk comptabilise la marge au moment de la vente de la note et couvre ensuite le risque de celle-ci~; ce que le client gagne ou
perd par la suite dépend du marché, que le \omterm{def:in:structurer-sales-and-sales-trader:structurer}{structureur} ne contrôle pas. Le lien avec le résultat du client passe en revanche
par la \omterm{def:in:structurer-sales-and-sales-trader:governance}{gouvernance des produits}, la réputation du desk et les ventes à venir.
\end{solution}

\begin{solution}{exo:in:structurer-sales-and-sales-trader:7}
Sur les 20 graines, la marge moyenne est de 1.84 pour 100, avec un écart-type de 0.06 (de 1.74 à 1.97)~: son signe
ne fait aucun doute, mais la deuxième décimale est incertaine.
\end{solution}

\begin{solution}{exo:in:structurer-sales-and-sales-trader:8}
La protection ne joue que si l'indice finit au-dessus de 60~\%~; en dessous, le porteur supporte dès le départ toute la baisse, si bien
que la note cumule les pires scénarios de l'indice sans ses meilleurs (elle est rappelée quand l'indice monte). Le coupon n'est
pas garanti non plus (il exige que l'indice soit au-dessus de 70~\%), et le porteur renonce aux dividendes. Il faut comparer la probabilité de perte et
la distribution, et non le coupon affiché.
\end{solution}

\begin{solution}{pb:in:structurer-sales-and-sales-trader:1}
\begin{enumerate}
\item Comme dans les définitions du chapitre.
\item Le client ou le distributeur, le desk de trading qui couvre la note, et la marge de la banque.
\item Exécute les ordres des clients sur le marché et en rend compte par rapport à leur indice de référence~; la gestion des ordres avec une personne dans la boucle
(livre~10, chapitre~20).
\item Valorisation et sensibilités~; traduction des anticipations et des produits~; mesure de l'exécution.
\item Livre~5, chapitre~18.
\item Cinq ans, observations annuelles, remboursement anticipé à 100~\%, coupon de 6~\% à 70~\% avec mémoire, capital perdu sous 60~\% à
l'échéance.
\item $100-V(c)$~; $c_m=(100-m-V(0))/(\partial V/\partial c)$ car la valeur est linéaire par rapport au coupon.
\item $-1.65$, 1.83, 6.91 et 9.13 pour 100.
\item 4.23~\%, 5.91~\%, 8.79~\% et 10.22~\%.
\item L'investisseur vend le risque de baisse et est rémunéré pour cela~; la probabilité de perdre son capital augmente aussi.
\item 486\,635 millions d'euros sur six marchés~; Suisse 58.8~\%, Allemagne 20.5~\%.
\item 2 à 10 millions d'euros par an, à une marge de 2~\%, sur 1 à 5 milliards d'euros d'émissions, avec dix \omterm{def:in:structurer-sales-and-sales-trader:structurer}{structureurs}.
\item Articles 16(3) et 24(2) (approbation et marché cible) et 25(1) (connaissances et compétence), avec les
lignes directrices de l'ESMA de 2023.
\item La Series 7 qualifie un représentant général en valeurs mobilières, avec la SIE~; la Series 63 couvre le droit des valeurs mobilières des États.
\item Salaire médian des agents commerciaux~: 103\,030~\$ dans le secteur des valeurs mobilières (90\up{e}~centile~: 309\,440~\$), 61\,440~\$ dans les banques.
\item À 6~\%~: $-1.65$ à 1~\% et 1.83 à 4~\% (20~\% de volatilité)~; 6.91 et 9.13 à 30~\%~; 2 à 10 millions d'euros par
\omterm{def:in:structurer-sales-and-sales-trader:structurer}{structureur}.
\item Il augmente la marge à coupon donné, ou le coupon à marge donnée, et il est aussi payé par l'acheteur.
\item Un produit adapté au marché cible, une description de ses risques que le client peut comprendre, et un prix juste (fair price).
\item La valorisation, le risque et la sensibilité aux paramètres que le client ne voit pas~; pas les relations et la distribution qui
déterminent le volume.
\item Le revenu du desk est la marge multipliée par le volume, et la marge dépend des taux et de la volatilité autant que de la conception.
Il faut accepter le poste si l'on veut travailler avec les clients et les produits~; les modèles comptent surtout pour dire
ce qu'une note vend réellement.
\end{enumerate}
\end{solution}

\begin{solution}{iq:in:structurer-sales-and-sales-trader:1}
Une obligation zéro-coupon qui verse 100 à l'échéance, plus une option d'achat sur l'indice (participation), achetée avec la différence
entre 100 et le prix de l'obligation~; la marge de l'émetteur est prélevée sur cette différence.
\iqlookfor{obligation plus option et budget.}
\end{solution}

\begin{solution}{iq:in:structurer-sales-and-sales-trader:2}
Des taux plus élevés (une part obligataire moins chère), une volatilité implicite plus élevée (la baisse vendue vaut davantage), et une conception différente
(une barrière plus basse ou un worst-of)~; la marge peut aussi varier.
\iqlookfor{taux, volatilité et conception.}
\end{solution}

\begin{solution}{iq:in:structurer-sales-and-sales-trader:3}
Vendeur de volatilité (vendeur du put sur la baisse)~; dans un worst-of, vendeur de corrélation~: une corrélation plus faible aggrave la pire performance
et rend la note moins chère, si bien que le coupon est plus élevé.
\iqlookfor{le signe de chaque exposition et pourquoi.}
\end{solution}

\begin{solution}{iq:in:structurer-sales-and-sales-trader:4}
Il faut comparer le mouvement du marché pendant l'exécution et l'impact attendu pour sa taille~: dérive défavorable, fourchette payée,
l'impact d'une participation de 5~\%~; puis montrer si l'algorithme s'est comporté comme prévu, avec l'analyse des coûts de transaction (TCA).
\iqlookfor{décomposition du slippage en mouvement de marché, fourchette et impact.}
\end{solution}

\begin{solution}{iq:in:structurer-sales-and-sales-trader:5}
Type de client, connaissances et expérience, capacité à supporter une perte (jusqu'à la totalité du capital sous 60~\%), tolérance au risque,
objectifs (revenu avec protection conditionnelle), horizon (jusqu'à cinq ans), et le marché cible négatif.
\iqlookfor{les catégories de marché cible MiFID appliquées à la note.}
\end{solution}

\begin{solution}{iq:in:structurer-sales-and-sales-trader:6}
Le desk est long de volatilité près des barrières et long de la baisse vendue par les investisseurs~; il détient un risque concentré de barrière
et de dividende. Il souffre lorsque l'indice baisse vers les barrières ensemble, lorsque le gamma bascule et que la couverture devient
coûteuse (livre~5, chapitre~27).
\iqlookfor{concentration du risque de barrière et position du desk dans la transaction.}
\end{solution}
